\section{Algorithms and Contracts}
\label{sec:algorithm}

This section gives the per-language outline pipeline, the significance
rules that classify nested constructs, the normative token-budget
truncation precedence, the sanitisation pass over verbatim file content,
the plan/execute pair semantics, and the compile-gate verdict shape.
The truncation precedence and sanitisation rules in particular are
\emph{normative}: two implementations must produce identical behaviour
on the same input. This is what allows \texttt{gt outline} run by the
CLI to match what the runtime hook prepends to the agent's
\texttt{Read} result, which in turn matches what
\texttt{codeindex\_explore} surfaces over MCP.

\subsection{Outline emission pipeline}
\label{sec:algorithm:pipeline}

\begin{algorithm}[H]
\KwIn{Source file path $\file$, target token budget $\budget$
  (default $300$).}
\KwOut{Markdown outline of $\file$ with $|\outline(\file)| \leq \budget$.}
\If{$\loc(\file) < \threshold$ \textbf{or} skip-comment present}{
  \Return $\emptyset$ \tcp*{outline not required}
}
$\lang \leftarrow \texttt{detectLang}(\file)$\;
$\textsc{ast} \leftarrow \texttt{parse}_\lang(\file)$\;
\If{$\textsc{ast}$ \textbf{is} \texttt{nil}}{
  \Return passthrough(\file) \tcp*{graceful degradation}
}
$\textsc{decls} \leftarrow \texttt{topLevelDecls}(\textsc{ast})$\;
$\textsc{sig} \leftarrow \texttt{significantNested}_\lang(\textsc{ast})$\;
$\textsc{doc} \leftarrow \texttt{sanitise}(\texttt{packageDoc}(\textsc{ast}))$\;
$\outline(\file) \leftarrow \texttt{render}(\textsc{decls}, \textsc{sig}, \textsc{doc})$\;
\While{$|\outline(\file)| > \budget$}{
  $\outline(\file) \leftarrow \texttt{truncateStep}(\outline(\file))$ \tcp*{per
    \cref{sec:algorithm:truncation}}
}
\Return $\outline(\file)$\;
\caption{The outline emission pipeline. \texttt{render} produces
  Markdown matching the schema in \cref{app:schema}.}
\label{alg:outline}
\end{algorithm}

\subsection{Per-language significance rules}
\label{sec:algorithm:significance}

The classifier \texttt{significantNested}$_\lang$ identifies nested
constructs worth surfacing in the outline. Per
\cref{assn:significance}, the per-language rules are fixed against a
specific tree-sitter grammar version (recorded in \texttt{go.mod}
and the in-tree parser-registration table). Grammar bumps require
updating the per-language node-kind tables in lockstep with the
acceptance-fixture suite.

The full per-language tables are in \cref{app:significance}. The
hybrid scheme is: per-language fixed list of significant constructs,
\emph{plus} a universal fallback for unrecognised constructs that have
line-span $\geq 10$ or carry a name/label.

\subsection{Token-budget truncation precedence}
\label{sec:algorithm:truncation}

When a candidate outline exceeds the budget $\budget$, content is cut
in this order until it fits. This precedence is \emph{normative}; two
implementations must produce identical truncation on the same input.

\begin{enumerate}[leftmargin=*]
  \item \textbf{Nested significant constructs inside private
        functions.} Lowest signal-to-token; private internals are
        usually irrelevant to outline navigation.
  \item \textbf{Nested significant constructs inside public functions,
        longest-first.} Replace with the literal token
        \texttt{[N~nested constructs truncated]}.
  \item \textbf{Private top-level decls collapsed to count.}
        \texttt{(+~12 private helpers)} replaces the individual
        entries.
  \item \textbf{Package-doc lines beyond the first three.} Verbatim
        text is the prompt-injection surface --- keep it bounded.
  \item \textbf{Imports collapsed to count.} \texttt{14 imports}
        replaces the comma-separated list.
\end{enumerate}

\Cref{app:schema} shows a worked specimen on the file
\texttt{cmd/gt/outline/walker\_go.go} (1530~LoC) under the default
$\budget = 300$.

\subsection{Sanitisation of verbatim file content}
\label{sec:algorithm:sanitization}

The outline lens auto-prepends to a \texttt{Read}, so any verbatim
text it surfaces (package doc, top-of-file comments, struct-field
doc strings) lands in agent context unconditionally --- including for
adversarial vendored or dependency code (\cref{assn:adversarial}).
Sanitisation is the cheapest mitigation, applied as a normative
pre-pass before render:

\begin{enumerate}[leftmargin=*]
  \item \textbf{Truncate.} Verbatim text per top-level decl is capped
        at $240$ characters.
  \item \textbf{Escape.} Lines matching prompt-injection patterns
        (literal strings \texttt{IGNORE PREVIOUS}, \texttt{SYSTEM:},
        \texttt{ASSISTANT:}, \texttt{USER:},
        \texttt{<|im\_start|>}, \texttt{<|im\_end|>}, and
        model-control tokens enumerated in \cref{app:sanitisation})
        are prefixed with the literal \texttt{[sanitised]~}.
  \item \textbf{Collapse.} Consecutive whitespace runs longer than
        three newlines collapse to two, defeating layout-based
        prompt-injection tricks.
\end{enumerate}

The sanitisation is conservative-biased: false positives (legitimate
prose flagged as injection) are acceptable; false negatives (injection
text emitted verbatim) are not. The \texttt{[sanitised]~} marker is
left in the output so an attentive reviewer can spot a flagged line.

\subsection{Plan/execute pair contract}
\label{sec:algorithm:planexecute}

For every compound op $\intent \in \intentset$ exposed as a pair
$\langle \intent, \intent! \rangle$:

\begin{enumerate}[leftmargin=*]
  \item $\intent(\file)$ emits a Markdown plan with five sections:
        \textbf{Target}, \textbf{Scope}, \textbf{Diff hunks},
        \textbf{Predicted gate verdict}, \textbf{Plan token}. The
        plan token is a content-addressed hash $h(\intent, \file,
        \diff)$ that the execute variant uses to identify which
        plan it is committing.
  \item $\intent!(\file, h)$ recomputes $\diff(\intent, \file)$
        from current file content, verifies the plan token matches,
        submits to gate $\gate$, and writes iff $\gate(\diff, \file)
        = \text{accept}$. If the plan token does not match (file
        changed since the plan was emitted), the execute aborts and
        the caller must re-plan.
\end{enumerate}

\begin{remark}
The plan-token mechanism is the trade-off in
\cref{sec:system:planexecute} made concrete. Stateless plans cost the
caller a re-plan when the file changes; the alternative --- a
stateful plan with a server-side cookie --- would couple the agent
to a session and complicate multi-agent dispatch. We chose
statelessness on the principle that an LLM that has to re-plan
because the file changed is recovering from a real conflict, not a
spurious one.
\end{remark}

\subsection{The compile gate}
\label{sec:algorithm:gate}

\begin{algorithm}[H]
\KwIn{Candidate diff $\diff$ over file $\file$; language $\lang$ of
  $\file$.}
\KwOut{Verdict $\verdict \in \{\text{accept}, \text{reject}\}$.}
$\file' \leftarrow \texttt{applyInTempfs}(\diff, \file)$\;
$\verdict_{\text{syntax}} \leftarrow \texttt{nativeSyntaxCheck}_\lang(\file')$\;
\If{$\verdict_{\text{syntax}} = \text{fail}$}{
  \Return $\text{reject}$\;
}
\If{$\lang$ \textbf{has} compile-aware LSP}{
  $\verdict_{\text{lsp}} \leftarrow \texttt{lspDiagnose}_\lang(\file')$\;
  \If{$\verdict_{\text{lsp}}$ \textbf{has new errors}}{
    \Return $\text{reject}$\;
  }
}
\Return $\text{accept}$\;
\caption{The compile-gate verdict. The temp-fs apply ensures the gate
  never touches the working tree; the candidate is materialised in a
  read-only scratch directory and discarded after the verdict.}
\label{alg:gate}
\end{algorithm}

The gate is conservative-biased by construction:
\cref{def:false_negative} permits it to reject safe diffs (false
negatives) but forbids accepting unsafe diffs (false positives). In
practice, native syntax check is the floor; the LSP-diagnose step is
opt-in for typed languages and provides a stronger gate when
available. The empirical false-negative rate is
\tothink{measured in \cref{sec:eval:gate}}.
